\begin{explanationbox}
	\textbf{\textcolor{CExplanation}{一句话直觉}}

	弦长公式就是把两点间距离与直线斜率结合起来，使计算只需横纵坐标之一即可完成。

	\medskip
	\textbf{\textcolor{CExplanation}{核心拆解}}
	\begin{description}[style=nextline, leftmargin=2.8em, labelsep=0.5em, itemsep=0.45em]
		\item[\textbf{要点一（两点距离）：}] 任意两点距离为
			\[
				|AB| = \sqrt{(x_1-x_2)^2 + (y_1-y_2)^2}.
			\]
		\item[\textbf{要点二（斜率代入）：}] 由直线方程得 $y_1-y_2 = k(x_1-x_2)$，代入距离公式：
			\[
				\begin{aligned}
					|AB| &= \sqrt{(x_1-x_2)^2 + k^2(x_1-x_2)^2} \\
					&= \sqrt{1+k^2}\,|x_1-x_2|.
			\end{aligned}
			\]
			同理，若 $k\neq0$，由 $x_1-x_2 = \frac{1}{k}(y_1-y_2)$ 可得另一形式。
	\end{description}

	\medskip
	\textbf{\textcolor{CExplanation}{几何本质（勾股比例）}}

	将弦长看作直角三角形斜边：水平直角边为 $|x_1-x_2|$，竖直直角边为 $|k||x_1-x_2|$，由勾股定理即得 $\sqrt{1+k^2}\,|x_1-x_2|$。

	\medskip
	\textbf{\textcolor{CExplanation}{代数意义（对称化简）}}

	公式将弦长表达为仅含横坐标（或仅含纵坐标）的形式，在联立直线与圆锥曲线后，可直接通过韦达定理整体计算，避免分别求交点坐标的繁琐。

	\medskip
	\textbf{\textcolor{CExplanation}{考点价值（提效工具）}}
	\begin{itemize}[leftmargin=*, itemsep=0.5em, topsep=0.4em]
		\item \textbf{考法一（联立求弦长）：} 联立直线与圆锥曲线，利用韦达定理求出 $x_1+x_2$ 和 $x_1x_2$，再计算 $\sqrt{1+k^2}\sqrt{(x_1+x_2)^2-4x_1x_2}$。
		\item \textbf{考法二（纵坐标版本）：} 在抛物线或斜率较复杂时，用 $|AB|=\sqrt{1+\frac{1}{k^2}}\,|y_1-y_2|$，可简化运算。
	\end{itemize}

	\medskip
	\textbf{\textcolor{CExplanation}{顿悟点}}

	\textbf{两个公式本质相同，核心是 $|AB|=\sqrt{1+k^2}\,|x_1-x_2|$；使用纵坐标版本务必检查 $k\neq0$，否则分母无意义。}

	\medskip
	\textbf{\textcolor{CExplanation}{使用场景}}
	\begin{itemize}[leftmargin=*, itemsep=0.5em, topsep=0.4em]
		\item \textbf{场景一（联立方程计算）：} 联立后已知两交点横坐标的和与积，直接用横坐标弦长公式。
		\item \textbf{场景二（纵坐标对称）：} 若曲线方程用纵坐标更方便（如 $y^2=2px$），可用纵坐标弦长公式。
	\end{itemize}
\end{explanationbox}
